\documentclass[11pt]{amsart}
\usepackage[margin=1.15in]{geometry}
\usepackage{amsmath,amssymb,amsthm,mathtools}
\usepackage{booktabs}
\usepackage[colorlinks,linkcolor=black,citecolor=black,urlcolor=blue!50!black]{hyperref}
\usepackage{xcolor}
\usepackage{lmodern}

\newtheorem{theorem}{Theorem}
\newtheorem{lemma}[theorem]{Lemma}
\newtheorem{proposition}[theorem]{Proposition}
\newtheorem{corollary}[theorem]{Corollary}
\theoremstyle{remark}
\newtheorem{remark}[theorem]{Remark}

\newcommand{\cbrt}[1]{\bigl\lfloor\sqrt[3]{#1}\bigr\rfloor}
\newcommand{\lean}[1]{\texttt{#1}}

\title[Grau Ribas's conjecture for A113768]{The cube-root recurrence OEIS A113768:\\ a proof of Grau Ribas's asymptotic conjecture}
\date{November 2026}
\thanks{The Lean development, this note and its \LaTeX{} source are public at
\url{https://github.com/tpicou/publicMaths}; see the folder \href{https://github.com/tpicou/publicMaths/tree/main/research/grau-ribas-2024}{\texttt{research/grau-ribas-2024}}.}

\begin{document}

\begin{abstract}
Let $a_1=1$ and $a_{n+1}=a_n+\cbrt{a_n}$. We determine the sequence exactly: within the
range $[k^3,(k+1)^3)$ it is an arithmetic progression of step $k$ whose first term is
$k^3+2^{\lceil\log_2 k\rceil}-k$, entered at index
$N_k=1+\tfrac32(k-1)(k+2)+\lceil\log_2 k\rceil$. Consequences are the conjecture of
J.~M.~Grau Ribas that $a_n\sim(2n/3)^{3/2}$, the fact that the perfect cubes in the sequence
are exactly the powers of~$8$, and the sharper expansion
$a_n=(2n/3)^{3/2}-n-\tfrac12\sqrt{2n/3}\,\log_2 n+O(\sqrt n)$.
All of these results, the last two with explicit constants, have been machine-checked in
Lean~4 with mathlib.
\end{abstract}

\maketitle

\section{Statements}

Throughout, $\cbrt{m}$ is the integer cube root of $m\in\mathbb N$, i.e.\ the unique
$k\in\mathbb N$ with $k^3\le m<(k+1)^3$. The sequence A113768 of the OEIS is
\[
  a_1=1,\qquad a_{n+1}=a_n+\cbrt{a_n}\quad(n\ge1),
\]
beginning $1,2,3,4,5,6,7,8,10,12,\dots,28,31,34,\dots$. The OEIS entry records the conjecture,
attributed to Jos\'e Mar\'ia Grau Ribas (2024), that $a_n$ grows like $(2n/3)^{3/2}$.

\begin{theorem}[Grau Ribas's conjecture]\label{thm:main}
$\displaystyle\lim_{n\to\infty}\frac{a_n}{(2n/3)^{3/2}}=1$. More precisely,
$a_n/(2n/3)^{3/2}=1+O(n^{-1/2})$; explicitly,
$\bigl|a_n/(2n/3)^{3/2}-1\bigr|\le 1800/\sqrt n$ for every $n\ge1$.
\end{theorem}

\begin{theorem}[Cubes]\label{thm:cubes}
A natural number $m$ is a perfect cube and a term of the sequence if and only if
$m=8^j$ for some $j\ge0$.
\end{theorem}

\begin{theorem}[Second-order asymptotics]\label{thm:second}
$\displaystyle a_n=\Bigl(\frac{2n}{3}\Bigr)^{3/2}-n-\frac12\sqrt{\frac{2n}{3}}\,\log_2 n+O(\sqrt n);$
explicitly, the $O(\sqrt n)$ term is at most $320\sqrt n$ in absolute value for every $n\ge1$.
\end{theorem}

All three follow from an exact description of the sequence (Proposition~\ref{prop:blocks}
and Proposition~\ref{prop:entry}). The constants $1800$ and $320$ are far from optimal; they are
the ones established in the formal proof (Section~\ref{sec:lean}), which follows the argument
below but replaces Taylor's theorem by the identity in Remark~\ref{rem:alg}.

\section{Notation and two arithmetic lemmas}

For $k\ge1$ let $\ell(k)=\lceil\log_2 k\rceil$, the least $\ell\ge0$ with $2^\ell\ge k$. Put
\[
  r_k=2^{\ell(k)}-k,\qquad
  \varepsilon_k=\begin{cases}1,&k\text{ is a power of }2,\\0,&\text{otherwise},\end{cases}\qquad
  w_k=3k+3+\varepsilon_k,
\]
and define $N_1=1$, $N_{k+1}=N_k+w_k$. Note that $\varepsilon_k=1$ if and only if $r_k=0$.

\begin{lemma}\label{lem:offset}
For every $k\ge1$:
\begin{enumerate}
  \item $0\le r_k<k$;
  \item $\ell(k+1)=\ell(k)+\varepsilon_k$;
  \item $r_{k+1}=k-1$ if $r_k=0$, and $r_{k+1}=r_k-1$ if $r_k>0$. In particular $r_{k+1}\le k-1$.
\end{enumerate}
\end{lemma}

\begin{proof}
(1) $r_k\ge0$ by definition of $\ell(k)$. If $k=1$ then $r_1=0$. If $k\ge2$ then $\ell(k)\ge1$
and minimality gives $2^{\ell(k)-1}<k$, so $2^{\ell(k)}<2k$, i.e.\ $r_k<k$.

(2) Write $\ell=\ell(k)$. If $k=2^\ell$ then $2^\ell<k+1\le2^{\ell+1}$, so $\ell(k+1)=\ell+1$.
Otherwise $2^\ell>k$, hence $2^\ell\ge k+1$; also $\ell\ge1$ (as $k=1$ is a power of $2$) and
$2^{\ell-1}<k<k+1$, so $\ell(k+1)=\ell$.

(3) If $r_k=0$ then $k=2^{\ell(k)}$ and by (2) $r_{k+1}=2k-(k+1)=k-1$. If $r_k>0$ then
$\ell(k+1)=\ell(k)$ and $r_{k+1}=2^{\ell(k)}-k-1=r_k-1$. In both cases $r_{k+1}\le k-1$ by (1).
\end{proof}

\begin{lemma}[Boundary identity]\label{lem:boundary}
For every $k\ge1$, \ $k^3+r_k+k\,w_k=(k+1)^3+r_{k+1}$.
\end{lemma}

\begin{proof}
If $r_k>0$ then $w_k=3k+3$ and the left side is $k^3+3k^2+3k+r_k=(k+1)^3+(r_k-1)$.
If $r_k=0$ then $w_k=3k+4$ and the left side is $k^3+3k^2+4k=(k+1)^3+(k-1)$.
Conclude with Lemma~\ref{lem:offset}(3).
\end{proof}

\section{The exact structure of the sequence}

Since $a_n\ge1$ implies $\cbrt{a_n}\ge1$, the sequence is strictly increasing and $a_n\ge n$.

\begin{proposition}[Blocks]\label{prop:blocks}
For every $k\ge1$ and every $0\le j\le w_k$,
\[
  a_{N_k+j}=k^3+r_k+jk,
\]
and $\cbrt{a_{N_k+j}}=k$ for $0\le j<w_k$. In particular $a_{N_k}=k^3+r_k$.
\end{proposition}

\begin{proof}
Induction on $k$. For $k=1$: $N_1=1$, $a_1=1=1^3+r_1$.

Suppose $a_{N_k}=k^3+r_k$ for some $k\ge1$; we show $a_{N_k+j}=k^3+r_k+jk$ for
$0\le j\le w_k$ by induction on $j$. The case $j=0$ is the hypothesis. If $j<w_k$ and the
formula holds for $j$, put $v=k^3+r_k+jk$. Then $v\ge k^3$, and by Lemma~\ref{lem:boundary}
and Lemma~\ref{lem:offset}(3),
\[
  v\le k^3+r_k+(w_k-1)k=(k+1)^3+r_{k+1}-k<(k+1)^3 .
\]
Hence $\cbrt v=k$, so $a_{N_k+j+1}=v+k=k^3+r_k+(j+1)k$. This also proves the cube-root claim.
Taking $j=w_k$ and using Lemma~\ref{lem:boundary},
$a_{N_{k+1}}=k^3+r_k+kw_k=(k+1)^3+r_{k+1}$, which is the hypothesis for $k+1$.
\end{proof}

\begin{corollary}\label{cor:blocks}
Since $N_1=1$ and $N_{k+1}>N_k$, every $n\ge1$ lies in exactly one interval $[N_k,N_{k+1})$,
and then $\cbrt{a_n}=k$, i.e.\ $k^3\le a_n<(k+1)^3$. For $k\ge2$ one has $a_{N_k-1}<k^3\le a_{N_k}$,
so $N_k$ is the first index at which the sequence reaches $k^3$.
\end{corollary}

\begin{proof}
The first two claims are Proposition~\ref{prop:blocks}. For the last, apply the proposition to
block $k-1$ with $j=w_{k-1}-1$: $a_{N_k-1}=k^3+r_k-(k-1)<k^3$ because $r_k\le k-2$ by
Lemma~\ref{lem:offset}(3).
\end{proof}

\begin{proposition}[Entry index]\label{prop:entry}
For every $k\ge1$,
\[
  N_k=1+\frac{3(k-1)(k+2)}{2}+\ell(k),\qquad\text{equivalently}\qquad 2N_k+4=3k^2+3k+2\ell(k).
\]
\end{proposition}

\begin{proof}
$N_k=1+\sum_{i=1}^{k-1}(3i+3)+\sum_{i=1}^{k-1}\varepsilon_i$. The first sum is
$\tfrac32(k-1)k+3(k-1)=\tfrac32(k-1)(k+2)$. The second counts the $j\ge0$ with $2^j\le k-1$,
i.e.\ with $2^j<k$, i.e.\ with $j<\ell(k)$; there are $\ell(k)$ of them.
\end{proof}

\begin{table}[ht]
\centering
\begin{tabular}{rrrrr}
\toprule
$k$ & $N_k$ & $a_{N_k}$ & $r_k$ & $w_k$\\
\midrule
1 & 1 & 1 & 0 & 7\\
2 & 8 & 8 & 0 & 10\\
3 & 18 & 28 & 1 & 12\\
4 & 30 & 64 & 0 & 16\\
5 & 46 & 128 & 3 & 18\\
6 & 64 & 218 & 2 & 21\\
7 & 85 & 344 & 1 & 24\\
8 & 109 & 512 & 0 & 28\\
9 & 137 & 736 & 7 & 30\\
10 & 167 & 1006 & 6 & 33\\
\bottomrule
\end{tabular}
\caption{The first ten blocks. The overshoot $r_k$ vanishes exactly at powers of two.}
\end{table}

\section{Proofs of the theorems}

\begin{proof}[Proof of Theorem~\ref{thm:cubes}]
Let $m=k^3$ be a term, $m=a_n$. As $a_n\ge1$, $k\ge1$, and $\cbrt{a_n}=k$, so
$N_k\le n<N_{k+1}$ by Corollary~\ref{cor:blocks}. Monotonicity and Proposition~\ref{prop:blocks}
give $k^3=a_n\ge a_{N_k}=k^3+r_k$, so $r_k=0$ and $k=2^j$, $m=8^j$. Conversely
$a_{N_{2^j}}=(2^j)^3+r_{2^j}=8^j$.
\end{proof}

\begin{proof}[Proof of Theorem~\ref{thm:main}]
Let $n\ge1$ and let $k$ be its block, so $N_k\le n\le N_{k+1}-1$ and $k^3\le a_n<(k+1)^3$.
From Proposition~\ref{prop:entry}, $2n\ge 2N_k=3k^2+3k-4+2\ell(k)\ge3k^2-1$, and, since
$2^k\ge k+1$ gives $\ell(k+1)\le k$,
\[
  2n\le 2N_{k+1}-2=3k^2+9k+2\ell(k+1)\le3k^2+11k .
\]
With $x=2n/3$ this reads $k^2\bigl(1-\tfrac1{3k^2}\bigr)\le x\le k^2\bigl(1+\tfrac{11}{3k}\bigr)$, hence
\[
  \Bigl(1+\frac{11}{3k}\Bigr)^{-3/2}\le\frac{a_n}{x^{3/2}}\le
  \Bigl(1+\frac1k\Bigr)^{3}\Bigl(1-\frac1{3k^2}\Bigr)^{-3/2}.
\]
Both bounds are $1+O(1/k)$. Finally $k\ge\cbrt{n}\to\infty$ because $a_n\ge n$, and
$k=\sqrt{x}\,(1+O(1/k))\asymp\sqrt n$, so the ratio is $1+O(n^{-1/2})$.
\end{proof}

\begin{proof}[Proof of Theorem~\ref{thm:second}]
Let $n=N_k+j$ with $0\le j<w_k\le3k+4$, and $x=2n/3$. By Proposition~\ref{prop:entry},
\[
  x=k^2+k+e,\qquad e=\tfrac13\bigl(2\ell(k)-4+2j\bigr),\qquad |e|\le3k,
\]
using $\ell(k)\le k$ (the finitely many $n$ with $k<8$ are absorbed in the $O$-term). Put $\delta=1/k+e/k^2$, so $x=k^2(1+\delta)$ and $|\delta|\le4/k$.
Assume $k\ge8$, so $|\delta|\le\frac12$. For $f(t)=(1+t)^{3/2}$ we have
$0<f''(t)=\tfrac34(1+t)^{-1/2}\le\tfrac34\sqrt2$ on $[-\tfrac12,\tfrac12]$, so Taylor's theorem gives
$f(\delta)=1+\tfrac32\delta+R$ with $|R|\le0.54\,\delta^2\le 9/k^2$. Therefore
\[
  x^{3/2}=k^3+\tfrac32k^2+\tfrac32ke+O(k)=k^3+\tfrac32k^2+jk+k\,\ell(k)+O(k).
\]
Since $a_n=k^3+r_k+jk$ with $0\le r_k<k$,
\[
  a_n-x^{3/2}=-\tfrac32k^2-k\,\ell(k)+O(k).
\]
Next, $n=\tfrac32x=\tfrac32k^2+\tfrac32(k+e)$ with $|k+e|\le4k$, so $-\tfrac32k^2=-n+O(k)$.
Moreover $\ell(k)=\log_2k+O(1)$ and $\log_2 n=2\log_2k+\log_2\tfrac32+\log_2(1+\delta)=2\log_2k+O(1)$,
so $k\,\ell(k)=\tfrac12k\log_2n+O(k)$. Finally $|\sqrt{1+\delta}-1|\le|\delta|$ gives
$|\sqrt x-k|\le4$, hence $\tfrac12k\log_2n=\tfrac12\sqrt x\log_2n+O(\log n)$. Collecting terms, and
using $k=O(\sqrt n)$,
\[
  a_n=x^{3/2}-n-\tfrac12\sqrt x\,\log_2n+O(\sqrt n).\qedhere
\]
\end{proof}

\begin{remark}\label{rem:alg}
For a fully explicit (and machine-checkable) version of the Taylor step, put $y=\sqrt x$ and
$P=k^3+\tfrac32k\,s$ with $s=x-k^2$. Then
\[
  (y^3-P)(y^3+P)=x^3-P^2=\tfrac34k^2s^2+s^3,
\]
and $-\tfrac13\le s\le5k$, $y^3+P\ge k^3/2$ give $|y^3-P|\le290k$. Likewise $(y-k)(y+k)=s$ gives
$k-1\le y\le k+3$. The Lean proof combines these polynomial bounds with
$\log_2k\le\ell(k)<\log_2k+1$ and $k^2\le n\le9k^2$.
\end{remark}

\begin{remark}
The $O(\sqrt n)$ term genuinely oscillates: it contains $r_k$ and the fractional part of
$\log_2k$. Numerically, the remainder divided by $\sqrt n$ stays between about $2.72$ and $3.18$
for $10^3\le n\le10^8$. The first-order statement explains the slow convergence of the ratio in
Theorem~\ref{thm:main}: $a_n/(2n/3)^{3/2}\approx1-\tfrac32(2n/3)^{-1/2}$, e.g.\ $0.99815$ at
$n=10^6$.
\end{remark}

\section{Formal verification}\label{sec:lean}

The Lean development is publicly available in the repository \url{https://github.com/tpicou/publicMaths}; all file
names below are relative to it and link to the files there. The file
\href{https://github.com/tpicou/publicMaths/blob/main/Maths/GrauRibas/CubeRoot.lean}{\texttt{Maths/GrauRibas/CubeRoot.lean}} formalises the sequence with
mathlib's integer cube root \lean{Nat.nthRoot 3}, and proves that it coincides with
$\lfloor m^{1/3}\rfloor$ computed in $\mathbb R$, so the Lean recurrence is literally the OEIS one.
The correspondence with this note is:

\medskip
\begin{center}
\begin{tabular}{ll}
\toprule
Lean declaration (namespace \lean{Maths.GrauRibas}) & Content\\
\midrule
\lean{a\_one}, \lean{a\_succ}, \lean{nthRoot\_three\_eq\_floor} & the recurrence with the real cube root\\
\lean{entry\_value}, \lean{block\_value} & Proposition~\ref{prop:blocks}\\
\lean{cbrt\_on\_block}, \lean{previous\_lt} & Corollary~\ref{cor:blocks}\\
\lean{entry\_formula} & Proposition~\ref{prop:entry}\\
\lean{cubes\_are\_powers\_of\_eight} & Theorem~\ref{thm:cubes}\\
\lean{grauRibas\_asymptotic} & Theorem~\ref{thm:main} (the limit)\\
\lean{ratio\_bound} & Theorem~\ref{thm:main} (rate, constant $1800$)\\
\lean{second\_order\_bound} & Theorem~\ref{thm:second} (constant $320$)\\
\bottomrule
\end{tabular}
\end{center}
\medskip

\smallskip
{\sloppy
\noindent\textbf{Trusted statements.} The file \href{https://github.com/tpicou/publicMaths/blob/main/Maths/GrauRibas/Spec.lean}{\texttt{Maths/GrauRibas/Spec.lean}} restates
Theorems~\ref{thm:main}--\ref{thm:second} (\lean{grau\_ribas\_conjecture}, \lean{rate},
\lean{cubes}, \lean{second\_order}) for an \emph{arbitrary} $f:\mathbb N\to\mathbb N$ with $f(1)=1$
and $f(n+1)=f(n)+\lfloor f(n)^{1/3}\rfloor$ ($n\ge1$), the cube root taken in $\mathbb R$. These
statements mention no auxiliary definition of the development, so checking them requires reading
only mathlib's definitions of floor, real powers, square root, $\log_2$ and limits;
\lean{exists\_solution} shows that the hypotheses are satisfiable.\par}

\smallskip
\noindent\textbf{Audit.} The project's Lean files contain no \lean{sorry}, \lean{admit},
\lean{native\_decide}, \lean{axiom}, \lean{unsafe}, \lean{implemented\_by}, \lean{opaque} or
\lean{set\_option} (in particular the default heartbeat limit is kept), and \lean{autoImplicit}
is disabled. The development builds with Lean~4.34.1 against the mathlib revision pinned in
\href{https://github.com/tpicou/publicMaths/blob/main/lake-manifest.json}{\texttt{lake-manifest.json}}; all listed theorems depend only on the standard axioms
\lean{propext}, \lean{Classical.choice} and \lean{Quot.sound}, and the whole development, together with every mathlib
declaration it uses, is accepted by an independent replay in the kernel
(\lean{leanchecker --fresh}). A kernel-checked computation also confirms the first 29 terms
against the OEIS data. To reproduce:
\begin{verbatim}
git clone https://github.com/tpicou/publicMaths
cd publicMaths
lake exe cache get
lake build
lake env leanchecker --fresh Maths
\end{verbatim}

\section*{Numerical check}
Independently of the proofs, the formulas for $N_k$ and $a_{N_k}$ were compared with a direct
computation of the first $3\cdot10^6$ terms (all $1412$ complete blocks agree), and the cubes
among them are exactly $8^0,\dots,8^{10}$. The same computation reproduces all $10\,000$ terms of
the OEIS b-file for A113768.

\begin{thebibliography}{9}
\bibitem{oeis} OEIS Foundation, \emph{The On-Line Encyclopedia of Integer Sequences}, sequence
A113768, \url{https://oeis.org/A113768}.
\bibitem{mathlib} The mathlib Community, \emph{The Lean mathematical library}, CPP 2020.
\end{thebibliography}

\end{document}
